\section{引言}

张量网络为高维数组的表示与计算提供了紧凑而有结构的形式，已经广泛用于量子多体计算、统计物理、图模型推断以及张量算法等方向\cite{orus2014practical}。
围绕张量网络收缩（TNC），已有工作主要关心三类问题：如何设计更好的收缩顺序，如何控制中间状态规模，以及如何通过中间状态压缩来降低近似收缩的代价\cite{markov2008simulating,dumitrescu2018benchmarking,jakes2019carving,gray2021hyperoptimized,huang2021efficient,pan2020contracting,gray2024hyperoptimized,ma2024approximate}。
现有研究大多围绕标量元素查询、局部可观测量计算或概率推断中的局部查询展开，而非显式生成保留开放腿后的完整输出张量\cite{orus2014practical,dechter1998bucket,pan2020contracting,heddes2026approximating}。
相比之下，一旦要求完整物化网络，计算与存储开销会显著膨胀\cite{jakes2019carving}。
本文关注的正是这一类带开放腿张量网络的完整输出计算问题。

这一问题与通常的近似收缩有一个根本差别：当完整输出本身就是任务目标时，算法不能将其替换为压缩表示，也不能仅返回少量局部结果。
因此，我们关注的问题不是改变输出形式，而是在保持完整输出要求不变的前提下，降低生成这一输出的计算开销。
若把端到端时间写成
\[
T_{\mathrm{total}} = T_{\mathrm{contract}} + T_{\mathrm{emit}},
\]
我们将 $T_{\mathrm{contract}}$ 理解为生成目标输出之前的内部计算开销，即网络内部的收缩、合并及相关中间状态处理；将 $T_{\mathrm{emit}}$ 理解为把这些结果显式物化为完整输出张量的开销。
由于完整输出本身必须被生成，$T_{\mathrm{emit}}$ 往往不可回避，因此方法设计的主要优化对象通常是 $T_{\mathrm{contract}}$。

基于这一观察，本文研究一种面向完整物化的近似收缩框架。
基本思路是：在网络交互图上构造一棵划分树，并沿树递归维护各个子网络的局部接口；在整个过程中，输出边界始终保持显式，只对内部接口做近似。
于是，本文近似的不是最终输出对象，而是“为了得到该输出而必须经过的内部中间状态”。
这也是本文与直接压缩输出端的方法，以及以标量收缩为目标的方法之间最本质的区别。

从表示上看，每个子树状态都可以理解为一个从输出侧到边界侧的接口矩阵；本文用低秩因子化在树上递归传播这一状态，并在每次合并后重新压缩，以抑制中间维度增长。
进一步地，针对内部节点合并时容易出现的维度膨胀，本文引入基于算子访问的隐式草图合并，不显式形成完整的合并接口，而是直接恢复其主导子空间。
因此，本文的核心并不是某一个孤立的局部技巧，而是把划分树递归、分隔状态近似与草图驱动的合并压缩组织成一条统一的完整收缩流程。
在这条流程里，输出语义保持不变，近似只服务于降低内部收缩成本。

实验部分围绕三个问题展开：第一，\method\ 在 TNC 任务中能否形成清楚而可用的精度--时间折中；第二，这些收益是否确实来自收缩阶段，以及哪些机制在较难网络上最关键；第三，近似强度是否能稳定控制误差。结果表明，在代表性经典网络上，\method\ 将相对误差控制在 $4.81\times 10^{-3}$ 以内，同时相对直接 TNC 获得 $13.98\times$ 到 $84.05\times$ 的加速；而固定秩基线采用固定小秩，虽然更快，却带来 $0.65$ 到 $0.98$ 的大误差。进一步分析表明，端到端收益几乎完全来自收缩阶段；隐式草图合并是稳定的效率来源，输出分块也具有实质作用。近似强度扫描实验则显示，近似强度是控制误差的旋钮。

本文的贡献可以概括为三点：
\begin{itemize}
    \item 提出一个面向带开放腿张量网络完整收缩的框架，将优化重点从“近似收缩值”转向“在保持输出边界语义不变的前提下降低内部 收缩代价”。
    \item 提出基于划分树递归的分隔状态框架，并在 \method\ 中引入基于算子访问的隐式草图合并，以降低内部合并阶段的大接口压缩开销；
    \item 通过主 benchmark、机制分析与强度扫描说明：\method\ 能够在 Ring-8、Tree-8、Grid $3\times 3$ 与 Random-8 上形成清楚的精度--时间折中，并揭示其收益来源与关键机制。
\end{itemize}
